\section*{Введение}
\addcontentsline{toc}{section}{Введение}

\textbf{Актуальность темы.} Рынок цифровой электроники характеризуется большим количеством
товарных позиций, быстрым обновлением моделей и высокой зависимостью покупательского выбора
от технических характеристик. Для магазина электроники важно не только представить каталог,
но и обеспечить своевременное обновление цен, остатков, категорий и статусов заказов.
Ручное ведение таких данных увеличивает вероятность ошибок, затрудняет обработку заявок и
снижает качество обслуживания покупателей. Поэтому разработка веб-приложения для управления
ассортиментом и заказами магазина цифровой электроники является актуальной практической задачей.

\textbf{Степень разработанности проблемы.} Современные интернет-магазины и маркетплейсы
предоставляют пользователям развитые средства поиска, фильтрации, оформления заказа и
отслеживания доставки. Однако готовые коммерческие платформы часто ограничивают владельца
магазина правилами конкретной площадки, комиссионной моделью или закрытой логикой управления.
Для небольшого магазина цифровой электроники актуальным остается создание самостоятельного
веб-приложения, которое учитывает особенности ассортимента, роли покупателей, менеджеров и
администраторов, а также требования к обработке заказов.

\textbf{Объект исследования} --- процессы управления онлайн-продажами цифровой электроники,
включая ведение ассортимента, взаимодействие покупателей с каталогом и обработку заказов.

\textbf{Предмет исследования} --- программные средства и архитектурные решения для разработки
веб-приложения, автоматизирующего управление товарами, категориями, пользователями, корзиной
и заказами магазина цифровой электроники.

\textbf{Целью данной работы} является проектирование и реализация веб-приложения для управления
ассортиментом и заказами магазина цифровой электроники, обеспечивающего покупателям удобную
работу с каталогом и заказами, а администраторам --- инструменты управления товарами,
категориями и заявками.
